\chapter{System Design and Pipeline Architecture}
\label{ch:design}

The pipeline consists of twelve numbered, single-purpose scripts in
\texttt{sepsis/project/}, preceded by the locked pre-registration document
(\texttt{00\_preregistration.md}). Each script reads from upstream outputs and
writes its own results to \texttt{sepsis/project/output/}. No script modifies the
original source data. All random seeds are fixed for reproducibility.
Figure~\ref{fig:pipeline_dag} shows the data flow, and
Appendix~\ref{app:stages} tabulates each stage's inputs and outputs in full.

This chapter describes what each stage does. It opens with why the architecture
takes this shape, because the structure is not incidental: a study whose
headline claim is that analyst decisions move results is obliged to make its own
analyst decisions inspectable, and several of the choices below exist to make a
specific class of silent error impossible rather than merely unlikely.


\section{Design Principles}
\label{sec:design_principles}

Four constraints shaped the architecture. Each is stated here with the decision
it forced and what that decision rules out.

\paragraph{The source data does not fit in memory, and must not be modified.}
MIMIC-IV's \texttt{chartevents} table is 3.5\,GB compressed and
\texttt{labevents} 2.6\,GB; expanded, either exceeds the working memory of the
machine this study ran on. Both are therefore read through a DuckDB streaming
layer that pushes the cohort filter and the item-identifier filter down into the
scan, so only rows belonging to the analysis cohort are ever materialised. The
alternative, extracting a reduced copy once and working from it, was
rejected because it makes the extraction step invisible to everything
downstream: a defect in the filter becomes indistinguishable from a property of
the data. Reading from the original tables on every run keeps the extraction
inside the reproducible chain.

The corollary is that the source directories are treated as strictly read-only.
No stage writes to \texttt{data/}, and every stage writes only to
\texttt{output/}. Reproduction is therefore a matter of deleting
\texttt{output/} and re-running, with no possibility that a previous run has
left the inputs in a state the next run depends on.

\paragraph{Every clinical constant must have exactly one definition.}
The study varies label definitions deliberately, which means the same quantity (an antibiotic--culture window, a SOFA threshold, an evaluation horizon)
is read by several stages and must mean the same thing in all of them. A
constant duplicated across two scripts is a constant that can silently diverge,
and in a study about definitional sensitivity such a divergence would be
indistinguishable from a finding. \texttt{config.R} is therefore the single
source of truth for every path, identifier, window, threshold, seed, family size
and event floor in the pipeline, and no other file hard-codes any of them. The
same discipline extends to the thesis: \texttt{thesis\_constants.R} writes the
reported numbers into the document from the result tables, so a value cannot
disagree between the code and the text either.

\paragraph{An interface between stages must fail loudly or not at all.}
The pipeline's most dangerous failure mode is not a crash but a stage that
quietly does less than it was asked to. Selecting predictors by intersecting a
declared feature list with the columns actually present, for instance, silently
drops any predictor whose name does not match: the model still fits, the
metrics still compute, and nothing in the output indicates that a variable is
missing. Interfaces are therefore declared rather than inferred:
\texttt{require\_features()} raises an error when a declared predictor is absent
instead of proceeding without it, and the fit diagnostics record convergence
status alongside the coefficient magnitudes that reveal an unidentified
likelihood, because a convergence flag alone does not distinguish the two.

The same principle governs the run-time invariants. Where a computation has a
property that must hold by construction, the pipeline asserts it and halts:
a threshold above every predicted probability must score exactly zero on the
Utility Score, because it \emph{is} the no-alert strategy; no operating point may
exceed the optimal strategy's score of one; and the three label variants must
produce three distinct sets of external onsets, since a label that was never
varied would otherwise present as a label that has no effect. These are stops,
not warnings. A warning in a sixty-minute batch run is a line in a log nobody
reads.

\paragraph{Expensive fits must be refreshable without disturbing the others.}
Fitting is the dominant cost, and gradient-boosted tree fitting is not
bit-stable across runs, so an incidental refit changes numbers throughout the
document without any change of substance. Two environment variables scope a run
accordingly: \texttt{V2\_VARIANTS} restricts the per-variant loops to a named
subset, and \texttt{V2\_ARMS} restricts the fitting arm
(Section~\ref{sec:ablation_method}). Adding an alternative label or an ablated
specification therefore does not re-execute the fits that every headline number
rests on, and the outputs of the additional arms are written to separate files
so they cannot overwrite a pre-registered result.

\paragraph{Why twelve scripts rather than one program.} The stage boundaries
follow the analysis, not the code: each corresponds to a step that has to be
separately inspectable because a claim in Chapter~\ref{ch:evaluation} rests on
it. The inference layer is the clearest case: Stage 12 reads only the stored
predictions of earlier stages, so intervals and multiplicity corrections can be
recomputed without refitting anything, and the correction applied to a result is
verifiable independently of the model that produced it. The cost of this
granularity is intermediate storage, which is paid in Parquet; the benefit is
that any single stage can be re-run and checked against the rest.

\begin{figure}[H]
    \centering
    \begin{tikzpicture}[
        node distance=0.4cm and 0.7cm,
        box/.style={rectangle, draw=ATUGreen, thick, rounded corners=2pt,
        fill=ATULightGreen!40, text width=2.6cm, align=center,
        font=\scriptsize, inner sep=3pt, minimum height=0.55cm},
        data/.style={rectangle, draw=ATUNavy, thick, rounded corners=2pt,
        fill=ATUTeal!20, text width=2.4cm, align=center,
        font=\tiny\itshape, inner sep=2pt, minimum height=0.45cm},
        arr/.style={-{Latex[length=1.5mm]}, thick, ATUGreen},
    ]

%--- Left column: scripts 01--05
        \node[data] (raw) {Raw data\\MIMIC-IV + eICU};
        \node[box, below=of raw]   (s01) {\texttt{01} Setup};
        \node[box, below=of s01]   (s02) {\texttt{02} Cohort Labels};
        \node[box, below=of s02]   (s03) {\texttt{03} Person-Hours};
        \node[box, below=of s03]   (s04) {\texttt{04} Sample Size};
        \node[box, below=of s04]   (s05) {\texttt{05} Primary Model};

%--- Right column: scripts 06--10
        \node[box, right=3.8cm of s01]  (s06) {\texttt{06} Comparators};
        \node[box, below=of s06]        (s07) {\texttt{07} Metrics Suite};
        \node[box, below=of s07]        (s08) {\texttt{08} External Valid.};
        \node[box, below=of s08]        (s09) {\texttt{09} Variance Decomp.};
        \node[box, below=of s09]        (s10) {\texttt{10} Equity Analysis};
        \node[box, below=of s10]        (s11) {\texttt{11} Clinical Utility};
        \node[box, below=of s11]        (s12) {\texttt{12} Inference};

%--- Output nodes (between columns)
        \node[data, right=of s02] (cohort)  {cohort\_\{A,B,C\}.pq};
        \node[data, right=of s03] (ph)      {person\_hours.pq};
        \node[data, right=of s05] (models)  {model\_primary.rds};
        \node[data, right=of s07] (metrics) {metric\_results.pq};

%--- Left-column flow
        \draw[arr] (raw)  -- (s01);
        \draw[arr] (s01)  -- (s02);
        \draw[arr] (s02)  -- (s03);
        \draw[arr] (s03)  -- (s04);
        \draw[arr] (s04)  -- (s05);

%--- Cross to right column
        \draw[arr] (s05)  -- (s06);

%--- Right-column flow
        \draw[arr] (s06)  -- (s07);
        \draw[arr] (s07)  -- (s08);
        \draw[arr] (s08)  -- (s09);
        \draw[arr] (s09)  -- (s10);
        \draw[arr] (s10)  -- (s11);
        \draw[arr] (s11)  -- (s12);

%--- Output arrows
        \draw[arr] (s02)  -- (cohort);
        \draw[arr] (s03)  -- (ph);
        \draw[arr] (s05)  -- (models);
        \draw[arr] (s07)  -- (metrics);

%--- Feed-forward: person-hours feed models and comparators
        \draw[arr, ATUNavy, densely dashed] (ph.east) to[out=0,in=180] (s06.west);

    \end{tikzpicture}
    \caption{Pipeline data flow. Each box is a single script; italic nodes are
    output files. Raw MIMIC-IV and eICU data are read-only.}
    \label{fig:pipeline_dag}
\end{figure}


\section{Configuration and Utilities}

\texttt{config.R} is the single source of truth for all pipeline settings:
data paths, cohort criteria, antibiotic name lists, MIMIC-IV item identifiers,
label variant definitions, split boundaries, prediction framing constants,
equity variable names, and Utility Score parameters. No other script hard-codes
a clinical constant or file path.

\texttt{utils.R} provides shared helper functions: a DuckDB-backed streaming
reader for large CSVs, Parquet reader/writer wrappers, datetime parsing,
last-observation-carried-forward (LOCF) imputation, and Schoenfeld test
wrappers.

The full pipeline can be reproduced on any MIMIC-IV installation by setting a
single environment variable (\texttt{MIMIC\_RESEARCH\_ROOT}).


\section{Stage 01: Setup Validation}

\texttt{01\_setup.Rmd} checks that all required MIMIC-IV source tables exist
and reports their sizes. The pipeline stops if any table is missing.


\section{Stage 02: Cohort Extraction and Labelling}

\texttt{02\_cohort\_labels.Rmd} builds all cohorts in a single pass:

\begin{enumerate}
    \item \textbf{Base cohort:} Adult first ICU stays with LOS $\geq 4$\,h
    (\pcNStays{} stays). Equity variables (race, language, insurance) are attached.

    \item \textbf{Suspected infection:} Qualifying antibiotic-culture pairs
    within each variant's time windows. Surveillance cultures are excluded.

    \item \textbf{SOFA organ dysfunction:} SOFA components are computed from
    chartevents, labevents, inputevents, and outputevents. A SOFA rise
    $\geq 2$ within each variant's windows defines organ dysfunction.

    \item \textbf{Onset time:} For the pre-registered variants A--C, the
    suspicion time $t_{\text{susp}}$, provided the SOFA rise occurs somewhere in
    the evaluation window. The SOFA criterion therefore gates \emph{membership}
    and contributes no timing information: an implementation choice with
    consequences developed in Section~\ref{sec:variant_f}.

    \item \textbf{Alternative labels (Protocol Amendments 2 and 4):} The Sepsis-3
    conjunction is also split into its two limbs
    (Section~\ref{sec:alt_labels}). Variant E keeps the suspicion anchor and
    drops the SOFA criterion; Variant D keeps organ dysfunction and drops the
    anchor, so its onsets are derived from physiology in Stage 03 rather than
    here. Variant F (Amendment 4) keeps the conjunction intact and re-times the
    onset to $\min(t_{\text{susp}}, t_{\text{SOFA}})$, also in Stage 03 because
    it needs the hourly SOFA trajectory. All three are held outside the
    pre-registered variant set in \texttt{config.R} and their outputs are
    written to separate files so they cannot enter H1 or the decomposition
    table.
\end{enumerate}

A single environment variable, \texttt{V2\_VARIANTS}, restricts the per-variant
loops in Stages 02--07 to a named subset. Adding the alternative labels
therefore did not require refitting variants A--C: XGBoost is not bit-stable
across runs, so an incidental refit would have moved every GBT-dependent number
in this thesis without any change of substance.


\section{Stage 03: Person-Hour Table Construction}

\texttt{03\_person\_hours.Rmd} expands each stay into hourly rows. Large source
tables (chartevents: 3.5\,GB; labevents: 2.6\,GB) are processed via DuckDB
streaming, pre-filtered to cohort stays, to avoid loading everything into
memory. For each person-hour:

\begin{itemize}
    \item Vitals and labs are LOCF-imputed from the most recent measurement.
    \item Six-hour rolling slopes are computed for respiratory rate, MAP, and
    lactate.
    \item Vasopressor indicators are set from concurrent infusions.
    \item Missingness flags are set for each lab test.
    \item The competing-event code (0/1/2/3) is assigned based on the hour's
    relationship to onset, discharge, or death.
    \item The temporal split group is assigned from the stay's admission year.
    \item Treatment anchor timestamps (first antibiotic, first culture) are
    recorded.
\end{itemize}

This is the pipeline's most expensive stage after model fitting, and the only
one that reads the full-size source tables; it is also the stage whose outputs
every later stage depends on, which is why the streaming and read-only
disciplines of Section~\ref{sec:design_principles} matter most here.


\section{Stage 04: Sample Size Calculation}

\texttt{04\_sample\_size.Rmd} applies the Riley et al.\
\cite{riley2019sample} criteria to the training portion (2008--2016) of each
label variant. The criteria assume independent observations, so the calculation
is parameterised on the \emph{stay} rather than the person-hour: the outcome
proportion supplied is the fraction of training stays labelled septic, and the
resulting minimum is compared against the number of training stays. The
per-person-hour parameterisation is computed alongside and retained as a
labelled reference, so that the two cannot be substituted for one another.


\section{Stage 05: Primary Model}

\texttt{05\_primary\_model.Rmd} fits the multinomial discrete-time
competing-risks model (Equation~\ref{eq:primary_model}) for each label variant
on the temporal training split. Baseline hazard splines use 5 knots, placed at
quantiles of the training hour distribution. A random patient-level 75/25 split
is also fitted to support H4.

The fit carries two guards that are part of the specification rather than
defensive extras. The declared predictor set is required rather than
intersected, so a name mismatch stops the stage instead of producing a model
with fewer terms than specified. And because the unpenalised likelihood for this
design is completely separated (its maximum lies at infinity, while the
optimiser's stopping rule reports success regardless
(Section~\ref{sec:primary_model_spec})) the fit is ridge-penalised and the
stage logs the largest fitted coefficient alongside the convergence flag, since
only the former distinguishes an identified solution from an arbitrary stopping
point. No analytic variance is available for the fitted object, so all
uncertainty for this model comes from the stay-level cluster bootstrap of
Stage 12.

Both stages 05 and 06 run two \textbf{arms}: the pre-registered specification
and the ablated one of Protocol Amendment 5
(Section~\ref{sec:ablation_method}), which drops the action-derived features.
The ablated arm writes to its own files and never overwrites a pre-registered
fit, so it can be added or refreshed without re-executing the fits every
headline number rests on.


\section{Stage 06: Comparator Models}

\texttt{06\_comparators.Rmd} fits the XGBoost GBT model, the static Cox model
(with Schoenfeld testing), and computes NEWS2, qSOFA, and SIRS at each
person-hour. GBT models are fitted under both temporal and random splits. GBT
feature importance is exported for each variant.

The gradient-boosted comparator is additionally fitted on the ablated feature
set with identical hyperparameters. NEWS2, qSOFA and SIRS never used the
ablated features, and the Cox comparator's role here is the
proportional-hazards check rather than discrimination, so neither is ablated.


\section{Stage 07: Metrics Suite}

\texttt{07\_metrics\_suite.Rmd} evaluates all models on all splits (temporal,
random, treatment-anchored) using the full metric set: AUROC, AUPRC, calibration
intercept and slope, PhysioNet Utility Score, alert burden, lead time, and
decision-curve net benefit. The ablated arm is scored by the same code on the
same test rows, and the paired difference is written to
\texttt{07\_ablation\_metrics.csv}.

The Utility Score is computed at the default 0.5 cut-off and swept across a
sixty-point grid of alert rates, in a single pass: the per-stay reward is
monotone in a falling threshold, so the whole grid costs one sort plus one
traversal of the reward-window rows. Two invariants are asserted at run time
rather than assumed: a threshold above every predicted probability must
return exactly zero, since it \emph{is} the no-alert strategy, and no operating
point may exceed the optimal strategy's score of one. A stand-alone Python
implementation of the same scoring rule, \texttt{utility\_score.py}, is kept in
step with the notebook so that prediction tables outside the pipeline can be
scored and the two implementations cross-checked.


\section{Stage 08: External Validation on eICU-CRD}

\texttt{08\_external\_validation.Rmd} applies the frozen MIMIC-IV models
(primary, GBT and NEWS2) to eICU-CRD without any retraining. It builds the
external person-hour panel once from physiology alone, then derives onsets
\emph{per label variant} from that panel using each variant's own
antibiotic--culture windows and SOFA-delta criterion
(Section~\ref{sec:label_transport}). Antibiotic and vasopressor exposure is
taken from \texttt{medication} and \texttt{infusionDrug} together.

Two arms are reported. The primary arm applies the Sepsis-3 conjunction to the
microbiology-covered sub-cohort; the secondary arm applies an antibiotic-only
anchor to the full cohort and is labelled as a different target. AUROC and
AUPRC are reported for each variant and model in both, alongside a
labelling-input coverage table and an explicit reconciliation of the eICU-CRD
and MIMIC-IV event rates.

The stage asserts at run time that the three variants produce distinct external
onset sets and halts otherwise, so that a label-invariant result can never be
produced by a label that was never varied.


\section{Stage 09: Variance Decomposition}

\texttt{09\_variance\_decomposition.Rmd} computes the pre-registered variance
decomposition by varying one factor at a time while holding all others fixed:

\begin{itemize}
    \item \textbf{Label spread (H1):} AUROC range across Variants A/B/C.
    \item \textbf{Model-class spread:} AUROC range across all models.
    \item \textbf{Split optimism (H4):} Random vs.\ temporal split comparison.
    \item \textbf{Metric divergence (H5):} AUROC vs.\ Utility Score decline.
\end{itemize}

It then runs the \textbf{label-anchor attribution} of Protocol Amendment 2:
performance under each limb of the Sepsis-3 conjunction, the
anchor-attributable fraction of Equation~\ref{eq:anchor_fraction}, and onset
agreement (Cohen's $\kappa$, median onset shift) between every label variant
and Variant B.


\section{Stage 10: Equity Analysis}

\texttt{10\_equity\_analysis.Rmd} evaluates model performance and label
sensitivity within each stratum of the equity variables declared in
\texttt{config.R} (race, ethnicity, language, insurance). MIMIC-IV v3.1 records
race and ethnicity in a single field, so three variables carry data and are
reported: \textbf{race, language, and insurance}. Levels are restricted to the
eight largest per variable, and are strata of an administrative field rather
than populations (Section~\ref{sec:equity_strata}). Discrimination is estimated
for every stratum with enough rows to support an estimate, but only strata
carrying at least \texttt{MIN\_SUBGROUP\_EVENTS\_INTERPRET} sepsis onsets enter
the exploratory contrast family, so the multiplicity correction is applied over
the analyses the study is willing to interpret rather than over every stratum
that happens to exist. Label sensitivity, being a stay-level proportion, is
reported for all strata.


\section{Stage 11: Clinical Utility Metrics}

\texttt{11\_clinical\_metrics.Rmd} computes flexible calibration curves,
decision-curve net benefit, operating points at a fixed per-hour sensitivity,
per-patient (event-level) detection and alarm-episode counts, and lead time
to onset, for both MIMIC-IV and eICU-CRD.


\section{Stage 12: Inference and Multiplicity}

\texttt{12\_inference.Rmd} is the inference layer. It reads only the Parquet
outputs of the earlier stages, so it can be re-run without refitting anything,
and it produces three things: stay-level cluster-bootstrap confidence intervals
for every headline and subgroup estimate; the Holm-corrected confirmatory tests
of H1--H6 and the Benjamini--Hochberg-corrected exploratory families (E1 and E2
subgroup analyses, E3 alternative-label contrasts); and the machine-readable
analysis register that supplies the status tags used
throughout Chapter~\ref{ch:evaluation}. The bootstrap resamples ICU stays, never
person-hours, and evaluates AUROC and AUPRC as weighted rank statistics on
pre-sorted data so that 2,000 replicates over half a million person-hours
remain tractable. Setting \texttt{V2\_BOOT\_B} overrides the replicate count
for a fast smoke test; reported results always use the values fixed in
\texttt{config.R}. The statistical rationale is given in
Section~\ref{sec:inference}.


\section{Reproducibility}

The pipeline is orchestrated by \texttt{run\_pipeline.R} (or the equivalent
\texttt{run\_pipeline.sh}), which executes all stages sequentially or by phase,
renders each stage to PDF in \texttt{output/pdf/}, and writes a per-stage and a
combined run log to \texttt{output/logs/}. All intermediate outputs are Parquet
(compressed, columnar). Small result tables ($<$1,000 rows) are additionally
written as CSV for readability. The end-to-end run reported in this thesis
completed in 60 minutes on an Apple silicon workstation; the dominant cost is
stage 05, which fits the multinomial model on 2.3 million person-hours for each
label variant.

\subsection{Fixed Random Seed}

All random operations use a single fixed seed (\texttt{42}) to ensure bitwise
reproducibility. Every Rmd script sets \texttt{set.seed(42)} in its setup
chunk before any computation. The same value is used for the patient-level
random split (\texttt{RANDOM\_SPLIT\_SEED} in \texttt{config.R}), the
\texttt{pmsampsize} sample-size calculation, and all model-fitting calls. No
seed is drawn from wall-clock time or left to the system default.

The full pipeline can be reproduced on any MIMIC-IV v3.1 installation by
setting the environment variable \texttt{MIMIC\_RESEARCH\_ROOT} to the
directory containing the \texttt{data/} folder.

\subsection{Code Availability}
\label{sec:code_availability}

The complete pipeline (the twelve numbered scripts, \texttt{config.R},
\texttt{utils.R}, \texttt{clinical\_scores.R}, \texttt{thesis\_constants.R},
\texttt{utility\_score.py}, the pre-registration document and the orchestration
scripts) is held in a public version-controlled repository at
\url{\repourl}, released under the \repolicence. Appendix~\ref{app:code}
describes the layout, the reproduction procedure, and the boundary of that
licence, which covers the source written for this study and neither the
databases nor the institutional template.

The repository contains \textbf{code only}. No MIMIC-IV or eICU-CRD record,
derived person-hour table, or fitted model object is included: the
\texttt{data/}, \texttt{processed\_data/} and \texttt{output/} directories are
excluded from version control by design, because the PhysioNet Credentialed
Data Use Agreement (Section~\ref{sec:ethics}) prohibits redistribution of the
underlying data. A reader wishing to reproduce these results must obtain
MIMIC-IV v3.1 and eICU-CRD v2.0 through their own PhysioNet credentialing, then
re-run the pipeline; with the fixed seed above it reproduces the reported
estimates exactly.
